\section{Background \& Related Work}
Naturally, the pursuit of structured generated outputs has two distinct branches
or "paradigms" to approach the problem: inference-time constraint decoding or
training-time weight alignment. Thus, we must analyze any computational
trade-offs in these approaches, and identify the critical latency-verification
gaps that necessitate a dual-objective framework like this to begin with.

% Review FSM based masking (outlines, JSON schemas, etc...)
  \subsection{Inference-Time Constraint Mechanisms}
  As aforementioned, the standard method for enforcing structural validity relies
  on projecting the generative vocabulary onto a formal grammar during the
  decoding step in the transformer. Frameworks such as Outlines~\cite{} and
  LMQL~\cite{} utilize Finite State Machines (FSMs) to compute a boolean mask over
  the vocabulary space $V$ at each timestep $t$. By mapping invalid token
  probabilities to $-\infty$, these systems can guarantee adherence to
  Context-Free Grammars (CFGs), such as standard JSON syntax, with negliglbe
  latency overhead (roughly $O(1)$ per token).

  However, the memoryless nature of FSM-based decoding renders it fundamentally
  incapable of enforcing Context-Sensitive Invariants (CSIs). Complex
  domain-specific languages (DSLs) and API schemas frequently demand cross-field
  arithmetic bounds, dynamic type resolution, or referential integrity (e.g.,
  referencing a dynamically instantiated variable in a subsequent payload block).
  By using an FSM, this means that these checks must be done \textit{post-hoc},
  where inadherence or malformed outputs would demand regeneration, and therefore
  significantly increase system latency. Additionally, attempting to model such
  state-dependent logic within regular languages like the aforementioned would of
  course result in a combinatorial state-space explosion, breaking the $O(1)$
  latency guarantee, and rendering standard FSM masking computationally impossible
  for any degree of richly typed semantics.

  Consequently, any viable formal verification framework must completely decouple 
  the enforcement of CFGs from CSIs, such as to avoid the combinatorial
  state-space explosion of modeling semantics within regular languages. 

  \subsection{Semantic-Aware Decoding}
  To enforce true operational semantics, recent literature has explored
  integrating target-side abstract machines directly into the generation loop.

  Scholak et al.~\cite{} introduced PICARD, which operates by parsing the top-$k$
  beams into Abstract Syntax Trees (ASTs) during decoding, and then actively
  bruning branches that violate database schemas. While this approach elevates the
  constraint boundary from syntax to semantics, it introduces a severe inference
  bottleneck. Because the core neural model remains oblivious to the underlying
  constraints, it still frequently proposes invalid distributions, and like
  mentioned above, can force the system into continuous, expensive rejection
  sampling.

  Similarly, Syncromesh~\cite{} employs Completion Engines (similar to above,
  performs constrained decoding, filtering and prefix validation) that query an
  external abstract machine to compute the exact admissible constraint space at
  every step. While mathematically sound, once again, the continuous token-level
  interaction between the Python-based neural model/inference loop and the
  symbolic state-tracker introduces significant overhead. Additionally, it lacks a
  decoding mechanisms that achieves the formal rigor of AST-verification without
  significantly degrading system throughput.

  To enforce these deferred semantics without falling into the latency traps of
  continuous token-level evaluation, significant optimization methods would be
  required. For instance, Kogler et al.~\cite{} decouples the verification
  process, by delegating deep semantic verification to discrete structural
  boundaries (element-level AST nodes), where if a completed element violates
  the CSI, the system triggers a speculative KV-cache rollback. This approach
  would allow the system to maintain the formal rigor of an abstract machine,
  while minimizing synchronization overhead.

  % How does LoRA work and is it really efficient?
  \subsection{Parameter-Efficient Fine-Tuning (PEFT)}
  As mentioned earlier, parallel to decoding research, efforts have also been
  focused on permanently shifting the model's predictive density ($P_\theta(Y|X)$)
  towards the target DSL.

  Patil et al~\cite{} demonstrated with Gorilla that Supervised Fine-Tuning (SFT)
  over API call traces can drastically reduce structural hallucination rates.
  Furthermore, Low-Rank Adaptation (LoRA))~\cite{} has been proven to be highly
  effective for this task. By freezing the pre-trained weights and injecting
  trainable rank decomposition matrices, LoRA allows models to internalize the
  rigid structural patterns of specific DSLs with minimal computational
  expenditure.

  However, purely neural approaches can suffer from stochastic decay---as even if  
  SFT could push a model to 99\% accuracy, the inherent probabilistic nature of
  transformers precludes truly formal guarantees. In safety-critical
  infrastructure, a 1\% failure rate constitutes a catastrophic system
  vulnerability. Attempting to close this gap using Reinforcement Learning (e.g.,
  CodeRL~\cite{}) has proven to be somewhat unstable, as strict operational
  semantics present a sparse, binary reward landscape that induces mode collapse
  during policy optimization. A potential solution might be to introduce richer
  end and transition states into training, such as to break out of the binary
  reward landscape.

  Ultimately, PEFT must be trated as an optimization mechanism rather than a
  standalone solution.

  % Approaches combining both?
  \subsection{Dual-Objective Architectures}
  Additionally, many advanced efforts to guarantee structural and semantic
  correctness attempt to unify neural predictive models with symbolic engines,
  effectively combining training-time alignment with inference-time verification,
  much more similar to what this paper sets out to do.

  AlphaGeometry~\cite{} couples a neural language model trained on millions of
  synthetic traces with a deterministic symbolic deduction engine. Similarly,
  frameworks such as Llema~\cite{} and Draft, Sketch and Prove~\cite{} fine-tune
  LLMs on highly structured mathematical datasets and interface them with formal
  interactive theorem provers, such as Lean or Coq.

  In these architectures, the neural model acts as a heuristic prior, navigating
  the combinatorial explosion of the search space, while the symbolic engine acts
  as the "oracle" of semantic truth. Furthermore, execution-guided synthesis
  frameworks like Jigsaw~\cite{} and Toolformer~\cite{} explicitly train the
  language model to anticipate the requirements of external verification APIs or
  compilers, thereby reducing the frequency of hallucinated, invalid proposals.

  Now, while these architectures successfully leveral neural training to aid
  symbolic verification, they consistently architect the interface between the two
  as an iterative, post-hoc communication loop rather than a deeply integrated
  decoding mechanism. In systems like the aforementioned AlphaGeometry or Llema,
  the LLM generates a complete candidate sequence, halts, and passes the payload
  to the symbolic engine for evaluation. Then, if the engine detects a semantic
  violation, it returns an error trace, and the LLM must condition on this error
  to regenerate~\cite{}.

  This macro-level propose-and-verify loop introduces three fatal flaws:
  \begin{itemize}
  \item Latency: The hand-offs between the GPU-bound neural model and CPU-bound
  symbolic verifier destroys inference speeds.
  \item Redundant compute: The LLM wastes massive computational cycles generate
  complete subtrees of tokens that might have semantic violations very early on,
  and therefore, the entire output is wrong.
  \item Lack of quantification: These frameworks fundamentally cannot formally
  measure how the neural model's policy alignment minimizes the symbolic engine's
  workload.
  \end{itemize}

  Another architecture, much closer to our approach, NeSyS~\cite{}, alternates
  training between a symbolic world model and a neural model, where the symbolic
  rules directly constrain the LLM's output probability distribution. The neural
  model is then fine-tuned specifically on trajectories inadequately covered by
  the rules. Similarly, Kogler et al.~\cite{} proposed a neuro-symbolic
  framework for DSL generation that integrates LLMs with a token-level syntactic
  validator and an element-level logical reasoning constraint solver.

  While frameworks like NeSyS demonstrate that symbolic engines can constrain
  neural outputs, their token-level probability modifications remain susceptible
  to the "Prefix Completeness Problem". Because the model only evaluates
  validity at the token boundary, it frequently misses multi-token semantic
  violations, trapping the transformer in unrecoverable KV-cache states.
  Conversely, Kogler et al.~\cite{} successfully bypasses this limitation by
  decoupling the verification process, using token-level syntactic validation
  alongside element-level logical reasoning.

  However, these frameworks still treat the neural model and constraint solver
  as static partners. They heavily optimize the inference engine to handle 
  constraint-solving but accept the computational overhead of the solver as a fixed, 
  unavoidable tax. They fail to leverage the neural policy to pro-actively
  bypass the verification bottleneck.
